\documentclass[11pt]{article}
\usepackage[a4paper,margin=1in]{geometry}
\usepackage{amsmath,amssymb,amsthm,booktabs,hyperref}
\hypersetup{colorlinks=true,linkcolor=blue,citecolor=blue,urlcolor=blue,pdftitle={The three-skeleton of the simplex on seven vertices is extendably shellable},pdfauthor={Xiaojian Zeng}}
\newtheorem{theorem}{Theorem}
\newtheorem{lemma}{Lemma}
\newtheorem{proposition}{Proposition}
\newcommand{\fac}{\mathcal F}
\title{The three-skeleton of the simplex on seven vertices is extendably shellable}
\author{Xiaojian Zeng\\PBC School of Finance, Tsinghua University, 100190\\\texttt{xiaojianzeng@foxmail.com}}
\date{}
\begin{document}
\maketitle
\begin{abstract}
We prove that the three-dimensional skeleton of the simplex on seven vertices is extendably shellable. This is the smallest parameter pair not covered by the previously known positive ranges for Simon's conjecture. The proof is a certified finite enumeration. We characterize a legal next facet by the inclusion-maximal intersections with the current facets, encode facet sets by 35-bit integers, and enumerate reachable sets modulo all 5040 vertex relabelings. There are 2,069,536 reachable orbits across the 35 layers; every proper reachable orbit has a legal successor. A 16 MB layer certificate and an independently implemented verifier accompany the result.
\end{abstract}

\section{Introduction}
Let $S_{d,n}$ be the $d$-skeleton of the $(n-1)$-simplex. Thus its facets are the $(d+1)$-subsets of $[n]$. Simon conjectured that every such skeleton is extendably shellable~\cite[Conjecture 4.2.1]{Simon1994}. The cases $d\leq2$ and $d\geq n-3$ are known; see \cite{Bigdeli2019,Culbertson2020,Coleman2022} and the references therein. Consequently $(n,d)=(7,3)$ is the first parameter pair outside those ranges.

The general conjecture was recently disproved by Bolognini and
Sentinelli: they exhibited a shellable subcomplex of $S_{3,16}$ whose
shelling cannot be extended, and proved more generally that $S_{d,d+13}$
is not extendably shellable for every $d\geq3$
\cite[Theorem~4.2 and Remark~4.4]{BologniniSentinelli2026}.  In arXiv
version~2, dated 31 July 2026, they identify $S_{3,7}$ as the first open
case.  We settle that case positively.

\begin{theorem}\label{thm:main}
Every shelling of a pure subcomplex generated by facets of $S_{3,7}$ extends to a shelling of $S_{3,7}$. Equivalently, $S_{3,7}$ is extendably shellable.
\end{theorem}

The argument is computer-assisted but uses only finite integer and set operations. No floating-point calculation, optimization, randomized sampling, or solver assertion enters the proof.

\section{The local extension criterion}
For a nonempty set $A\subseteq\fac(S_{3,7})$, write $\langle A\rangle$ for the subcomplex generated by $A$. If $F\notin A$, set
\[ I_A(F)=\{F\cap G:G\in A\}. \]

\begin{lemma}\label{lem:local}
The facet $F$ may be appended after a shelling whose facet set is $A$ if and only if every inclusion-maximal member of $I_A(F)$ has cardinality three. Equivalently, $I_A(F)$ contains a three-set and every member of $I_A(F)$ is contained in a three-set belonging to $I_A(F)$.
\end{lemma}
\begin{proof}
The intersection $\langle F\rangle\cap\langle A\rangle$ is generated inside the tetrahedron $F$ by the faces $F\cap G$, $G\in A$. Its facets are exactly the inclusion-maximal members of $I_A(F)$. The shelling condition says that this intersection is nonempty and pure of dimension two, which is precisely that all these maximal members have size three. The second formulation is the elementary characterization of a finite family whose maximal members all have size three.
\end{proof}
Thus whether a next facet is legal depends only on $A$, not on the shelling order that reached $A$.

\section{Certified orbit enumeration}
Order the 35 four-subsets of $[7]$ lexicographically and represent a facet set by the corresponding 35-bit unsigned integer. Let $L_k$ be the set of canonical representatives, under the natural action of $\mathfrak S_7$, of the $k$-facet sets reachable by a partial shelling. The canonical representative is the numerically least bit string in its orbit. Vertex transitivity gives $L_1=\{1\}$.

Starting with $L_1$, the enumeration forms every $A\cup\{F\}$ with $A\in L_k$ and $F$ legal by Lemma~\ref{lem:local}, canonicalizes the result under all 5040 permutations, and removes duplicates. This defines $L_{k+1}$ exactly. The largest layer is $L_{20}$, with 337,584 orbits. In total the layers contain 2,069,536 orbits, and the enumeration examines 29,319,999 legal transitions. Selected layer sizes are shown in Table~\ref{tab:layers}; the complete list accompanies the certificate.

\begin{table}[ht]\centering
\caption{Selected orbit-layer sizes.}\label{tab:layers}
\begin{tabular}{rrrrrrrr}\toprule
$k$&1&5&10&15&20&25&30\\
$|L_k|$&1&22&2,404&58,803&337,584&38,203&137\\\midrule
$k$&31&32&33&34&35&&\\
$|L_k|$&38&10&3&1&1&&\\\bottomrule
\end{tabular}\end{table}

\begin{proposition}\label{prop:certificate}
For every $1\leq k<35$ and every $A\in L_k$, at least one missing facet is legal. Moreover, the unique member of $L_{35}$ is the full facet set.
\end{proposition}
\begin{proof}
The certificate stores every $L_k$ as a strictly increasing sequence of little-endian unsigned 64-bit integers. A manifest records the byte length and digest of each layer file. The independent verifier first checks the manifest and strict uniqueness of each layer. It recomputes legality from the inclusion-maximal elements of $I_A(F)$, rather than using the ridge-cover implementation of the generator. For each $k<35$ it checks that no $A\in L_k$ has zero legal moves and that the canonicalized set of all recomputed successors is exactly the stored $L_{k+1}$. Finally it checks that $L_{35}$ consists of the integer with its lowest 35 bits all equal to one. All comparisons are exact. The verifier reports success for all 34 transitions.
\end{proof}

\begin{proof}[Proof of Theorem~\ref{thm:main}]
Let a partial shelling have facet set $A$ and size $k$. Relabeling vertices maps $A$ to its representative in $L_k$ and preserves Lemma~\ref{lem:local}. If $k<35$, Proposition~\ref{prop:certificate} gives a legal next facet. Append it and iterate. After $35-k$ steps the full skeleton is obtained.
\end{proof}

\section{Reproducibility and limitations}
The archive contains the generator, independent verifier, 35 binary layer files, manifest, raw output, and timing logs. From the archive root, the decisive command is
\begin{verbatim}
python3 certificate/verify_layers.py certificate/layers \
  --output certificate/verification.json
\end{verbatim}
The recorded verification used CPython 3.10.12, NumPy 2.2.6, and Numba 0.67.0 on x86-64 Linux. It completed in 6 minutes 29 seconds with peak resident set size 308 MB. A separate exhaustive cross-check verifies the equivalence of the two local criteria on all nonempty facet states for the five-vertex analogue.

The result is restricted to $S_{3,7}$; it does not restore the general
Simon conjecture, which is false~\cite{BologniniSentinelli2026}.  The
released exact certificate and independent verifier define the
computational proof and reproduce every finite check used here.

\bibliographystyle{plain}
\bibliography{references}
\end{document}
